\documentclass[11pt,a4paper]{article}
\usepackage[utf8]{inputenc}
\usepackage{amsmath,amssymb,amsthm}
\usepackage{listings}
\usepackage{xcolor}
\usepackage[margin=2.5cm]{geometry}
\usepackage{hyperref}

\lstset{
  language=C++,
  basicstyle=\ttfamily\small,
  keywordstyle=\color{blue}\bfseries,
  commentstyle=\color{green!60!black},
  stringstyle=\color{red!70!black},
  numbers=left,
  numberstyle=\tiny\color{gray},
  breaklines=true,
  frame=single,
  tabsize=4,
  showstringspaces=false
}

\title{IOI 1999 -- Road Network (Tree Median)}
\author{Solution}
\date{}

\begin{document}
\maketitle

\section{Problem Statement}

Given a weighted tree with $N$ nodes, find the node $v$ that minimizes the sum of
distances to all other nodes:
\[
  S(v) = \sum_{u=1}^{N} d(v, u).
\]
Output that node and the minimum sum.

\section{Solution Approach}

\subsection{Rerooting Technique}

A naive approach computes $S(v)$ independently for each $v$, taking $O(N^2)$ time.
The rerooting technique reduces this to $O(N)$.

\paragraph{Pass 1 -- Root at node 1.} Compute via DFS:
\begin{itemize}
  \item $\mathrm{sz}[v]$: number of nodes in the subtree rooted at $v$.
  \item $\mathrm{subSum}[v]$: sum of distances from $v$ to all nodes in its subtree.
\end{itemize}
Then $S(1) = \mathrm{subSum}[1]$.

\paragraph{Pass 2 -- Reroot.} When moving from parent $p$ to child $c$ along an
edge of weight $w$:
\[
  S(c) = S(p) + w \cdot (N - 2 \cdot \mathrm{sz}[c]).
\]

\begin{proof}[Derivation]
Relative to $p$, all $\mathrm{sz}[c]$ nodes in $c$'s subtree become $w$ closer
to $c$, contributing $-w \cdot \mathrm{sz}[c]$. The remaining $N - \mathrm{sz}[c]$
nodes become $w$ farther, contributing $+w \cdot (N - \mathrm{sz}[c])$. Combining:
\[
  S(c) = S(p) - w \cdot \mathrm{sz}[c] + w \cdot (N - \mathrm{sz}[c])
       = S(p) + w \cdot (N - 2\,\mathrm{sz}[c]).
\]
\end{proof}

\section{C++ Solution}

\begin{lstlisting}
#include <cstdio>
#include <cstring>
#include <vector>
#include <stack>
#include <algorithm>
using namespace std;

int main() {
    int N;
    scanf("%d", &N);

    vector<vector<pair<int, long long>>> adj(N + 1);
    for (int i = 0; i < N - 1; i++) {
        int u, v;
        long long w;
        scanf("%d %d %lld", &u, &v, &w);
        adj[u].push_back({v, w});
        adj[v].push_back({u, w});
    }

    vector<long long> sz(N + 1, 0), subSum(N + 1, 0), dist(N + 1, 0);
    vector<int> order, parent(N + 1, 0);
    vector<long long> parentW(N + 1, 0);
    vector<bool> visited(N + 1, false);

    // Pass 1: BFS/DFS from node 1 to get traversal order
    stack<int> stk;
    stk.push(1);
    visited[1] = true;
    while (!stk.empty()) {
        int u = stk.top(); stk.pop();
        order.push_back(u);
        for (auto& [v, w] : adj[u]) {
            if (!visited[v]) {
                visited[v] = true;
                parent[v] = u;
                parentW[v] = w;
                stk.push(v);
            }
        }
    }

    // Compute subtree sizes and sums (process leaves first)
    for (int i = (int)order.size() - 1; i >= 0; i--) {
        int u = order[i];
        sz[u] = 1;
        subSum[u] = 0;
        for (auto& [v, w] : adj[u]) {
            if (v != parent[u]) {
                sz[u] += sz[v];
                subSum[u] += subSum[v] + w * sz[v];
            }
        }
    }

    // Pass 2: reroot to compute dist[u] = S(u)
    dist[1] = subSum[1];
    for (int i = 0; i < (int)order.size(); i++) {
        int u = order[i];
        for (auto& [v, w] : adj[u]) {
            if (v != parent[u]) {
                dist[v] = dist[u] + w * (N - 2 * sz[v]);
            }
        }
    }

    // Find the node minimizing S(v)
    int bestNode = 1;
    long long bestDist = dist[1];
    for (int u = 1; u <= N; u++) {
        if (dist[u] < bestDist) {
            bestDist = dist[u];
            bestNode = u;
        }
    }

    printf("%d\n%lld\n", bestNode, bestDist);

    return 0;
}
\end{lstlisting}

\section{Complexity Analysis}

\begin{itemize}
  \item \textbf{Time complexity:} $O(N)$. Two linear passes over the tree.
  \item \textbf{Space complexity:} $O(N)$ for adjacency lists and auxiliary arrays.
\end{itemize}

\end{document}
